\chapter*{คำนำ}
\addcontentsline{toc}{chapter}{คำนำ}

วิชาฟิสิกส์พื้นฐานเป็นองค์ความรู้ที่เป็นรากฐานสำคัญของวิทยาศาสตร์ทุกแขนง โดยเฉพาะอย่างยิ่งสำหรับนักศึกษาสาขาวิทยาศาสตร์ชีวภาพ จุลชีววิทยา สิ่งแวดล้อม และเกษตรศาสตร์ การเข้าใจปรากฏการณ์ทางกายภาพของสิ่งมีชีวิตและสภาพแวดล้อม ไม่ว่าจะเป็นการไหลของของเหลวในท่อลำเลียงของพืช แรงดันออสโมซิสในเซลล์ การทำงานของเครื่องหมุนเหวี่ยงแยกอนุภาค การลำเลียงประจุและศักย์ไฟฟ้าที่เยื่อหุ้มเซลล์ กลไกการเกิดภาพในกล้องจุลทรรศน์แบบใช้แสง ไปจนถึงปฏิกิริยานิวเคลียร์และการแผ่รังสีที่ใช้ในการถนอมอาหารและการปรับปรุงพันธุ์พืช ล้วนต้องอาศัยกฎและหลักการทางฟิสิกส์ทั้งสิ้น

ตำราและคู่มือปฏิบัติการเรื่อง ``ฟิสิกส์พื้นฐานและคู่มือปฏิบัติการสำหรับวิทยาศาสตร์ชีวภาพและสิ่งแวดล้อม'' เล่มนี้ ได้รับการเรียบเรียงขึ้นอย่างเป็นระบบตามมาตรฐานวิชาการของมหาวิทยาลัยราชภัฏรำไพพรรณี โดยมุ่งเน้นการบูรณาการหลักการฟิสิกส์พื้นฐานเข้ากับการประยุกต์ใช้จริงในสิ่งมีชีวิตและกระบวนการทางชีวภาพ เนื้อหาครอบคลุมทั้งหมด 13 บทเรียนอย่างสมบูรณ์ ตั้งแต่ระบบหน่วยและการวัด เวกเตอร์ จลนศาสตร์ กฎการเคลื่อนที่ของนิวตัน โมเมนตัม สมดุลกล งานและพลังงาน กำลังและเครื่องกลอย่างง่าย ปรากฏการณ์ความร้อน อุณหพลศาสตร์ ไฟฟ้าสถิต ไฟฟ้ากระแส ไปจนถึงทัศนศาสตร์และกัมมันตภาพรังสี

จุดเด่นสำคัญของตำราเล่มนี้คือการผสานชุดการทดลองและคู่มือปฏิบัติการวัดละเอียด (Precision Measurement Laboratory Manual) ไว้อย่างเป็นรูปธรรม พร้อมตัวอย่างการคำนวณและข้อผิดพลาดที่พบบ่อย เพื่อให้นักศึกษาสามารถเชื่อมโยงทฤษฎีเชิงคณิตศาสตร์เข้ากับการทดลองจริงในห้องปฏิบัติการได้อย่างมีประสิทธิผล ผู้เขียนหวังเป็นอย่างยิ่งว่าตำราเล่มนี้จะเป็นประโยชน์สูงสุดแก่คณาจารย์ นักศึกษา และนักวิจัยในการพัฒนาองค์ความรู้ทางวิทยาศาสตร์สืบไป

\vspace{1.5cm}
\begin{flushright}
    \textbf{ผศ.ดร.ชีวะ ทัศนา}\\[0.2cm]
    สาขาวิชาฟิสิกส์ คณะวิทยาศาสตร์และเทคโนโลยี\\
    มหาวิทยาลัยราชภัฏรำไพพรรณี\\
    พ.ศ. 2569
\end{flushright}
